\documentclass[11pt,aspectratio=169]{beamer}

\usepackage{xcolor}
\definecolor{gradleteal}{HTML}{02303A}
\definecolor{gradlegreen}{HTML}{1FA87A}
\definecolor{codelight}{HTML}{F4F7F7}
\definecolor{codeframe}{HTML}{9FB8B8}
\definecolor{codecomment}{HTML}{4F4F4F}
\definecolor{codestring}{HTML}{A31515}
\definecolor{codekeyword}{HTML}{0033B3}
\definecolor{linkgreen}{HTML}{0B6E4F}

\usetheme{Madrid}
\usecolortheme[named=gradleteal]{structure}
\setbeamertemplate{navigation symbols}{}

\usepackage{listings}
\usepackage{hyperref}
\hypersetup{colorlinks=true,linkcolor=gradleteal,urlcolor=linkgreen,citecolor=gradleteal}

\setbeamercolor{title}{fg=gradleteal,bg=white}
\setbeamercolor{subtitle}{fg=gray,bg=white}
\setbeamercolor{author}{fg=black,bg=white}
\setbeamercolor{date}{fg=gray,bg=white}
\setbeamertemplate{title page}{
  \vfill
  \begin{center}
    {\usebeamerfont{title}\color{gradleteal}\inserttitle}\par
    \vskip2em
    {\usebeamerfont{author}\color{black}\insertauthor}\par
    \vskip1em
    {\usebeamerfont{date}\color{gray}\insertdate}
  \end{center}
  \vfill
}
\setbeamercolor{frametitle}{fg=gradleteal,bg=codelight}
\setbeamercolor{block title}{fg=white,bg=gradleteal}
\setbeamercolor{block body}{fg=black,bg=codelight}
\setbeamercolor{item}{fg=gradleteal}
\setbeamercolor{author in head/foot}{fg=white,bg=gradleteal}
\setbeamercolor{title in head/foot}{fg=white,bg=gradleteal}
\setbeamercolor{date in head/foot}{fg=white,bg=gradleteal}
\setbeamercolor{section in head/foot}{fg=white,bg=gradleteal}

\lstset{
  backgroundcolor=\color{codelight},
  rulecolor=\color{codeframe},
  basicstyle=\ttfamily\scriptsize,
  keywordstyle=\color{codekeyword}\bfseries,
  commentstyle=\color{codecomment}\itshape,
  stringstyle=\color{codestring},
  breaklines=true,
  frame=single,
  showstringspaces=false,
  tabsize=2,
}
\lstdefinestyle{shell}{
  language=bash,
  keywordstyle=\color{black},
  morekeywords={gradle,gradlew,java,javac,git,cd,mkdir,code,sudo,apt},
}

\title[Gradle in CS F213 Lab]{Using Gradle in CS F213 Lab}
\author[A. Pathak]{Aman Pathak, PhD TA \\ \href{mailto:p20260018@goa.bits-pilani.ac.in}{\texttt{p20260018@goa.bits-pilani.ac.in}}}
\date{CS F213 Lab}

\begin{document}

\begin{frame}[plain]
  \titlepage
\end{frame}

\begin{frame}[fragile]{The problem: building by hand does not scale}
  For one file, this is fine:
  \begin{lstlisting}
javac Calculator.java
java Calculator
  \end{lstlisting}
  Real projects break this habit quickly:
  \begin{itemize}
    \item many files in packages, compiled in the right order
    \item third-party libraries (where do the jars live?)
    \item tests that must run the same way on every machine
    \item ``it works on my laptop'' is not a build strategy
  \end{itemize}
  A build tool takes over all of this. Ours is Gradle, and we drive it
  from the terminal.
\end{frame}

\begin{frame}{Why Gradle: what you get}
  You describe \emph{what} the project is, Gradle figures out \emph{how}
  to build it:
  \begin{itemize}
    \item \textbf{One setup command:} \texttt{gradle init} scaffolds the
      whole project, folders and build files included
    \item \textbf{Dependencies by name:} libraries fetched from Maven
      Central, no jars passed around by hand
    \item \textbf{Standard tasks:} \texttt{build}, \texttt{test},
      \texttt{run}, \texttt{clean} work the same in every project
    \item \textbf{Incremental builds:} only what changed gets rebuilt
    \item \textbf{Same build everywhere:} the wrapper pins one Gradle
      version for the whole class
  \end{itemize}
\end{frame}

\begin{frame}[fragile]{gradle vs.\ gradlew: the one distinction to remember}
  \begin{center}
  {\scriptsize
  \begin{tabular}{l|l|l}
    & \texttt{gradle} & \texttt{./gradlew} \\ \hline
    What it is & system install & wrapper script inside the repo \\
    Version & whatever that machine has & pinned in
      \texttt{gradle/wrapper/*.properties} \\
    Comes from & pre-installed & created by \texttt{gradle init} \\
  \end{tabular}
  }
  \end{center}
  \vspace{0.5em}
  \textbf{Rule for this course:} set the project up with \texttt{gradle},
  then do everything else with \texttt{./gradlew}. Same command, same
  version, every machine:
  \begin{lstlisting}
gradle init       # once: scaffolds the project (needs system Gradle)
./gradlew build   # always: builds with the pinned version
  \end{lstlisting}
  Both commands must be run \textbf{inside the project root} (the folder
  that holds \texttt{settings.gradle.kts}).

  \vspace{0.3em}
  {\scriptsize Further reading: \href{https://docs.gradle.org/current/userguide/gradle_wrapper.html}{Gradle Wrapper}
  \textbar{} \href{https://docs.gradle.org/current/userguide/build_init_plugin.html}{Build Init plugin}}
\end{frame}

\begin{frame}[fragile]{Step 0: check your tools}
  \begin{lstlisting}
java -version      # we use Java 21
gradle --version  # we use Gradle 9.x
  \end{lstlisting}
  You need exactly two VS Code extensions, nothing more:
  \begin{itemize}
    \item Microsoft Extension Pack for Java (language support, debugger,
      test runner; Red Hat support is already inside it)
    \item Gradle for Java (the Gradle tasks panel)
  \end{itemize}
  Do not use Code Runner: it compiles single files directly and ignores
  the build files and the test framework.
\end{frame}

\begin{frame}[fragile]{Step 1: create the project with \texttt{gradle init}}
  Make an empty folder, step into it, and run one command:
  \begin{lstlisting}
mkdir gradle-demo
cd gradle-demo
gradle init
  \end{lstlisting}
  Gradle now asks a short series of questions. Answer them exactly like
  this (plain Enter means ``take the default''):
  \begin{center}
  {\scriptsize
  \begin{tabular}{l|l}
    Question & Your answer \\ \hline
    Type of build & 1: Application \\
    Implementation language & 1: Java \\
    Target Java version & Enter (default: 21) \\
    Project name & Enter (default: folder name) \\
    Application structure & 1: Single application project \\
    Build script DSL & 1: Kotlin (or Enter, it is the default) \\
    Test framework & 1: JUnit 4 \\
    New APIs and behavior & Enter (default: no) \\
  \end{tabular}
  }
  \end{center}
\end{frame}

\begin{frame}[fragile]{Step 2: look at what \texttt{init} just made}
  You wrote no files yourself. The project now looks like this:
  \begin{lstlisting}
gradle-demo/
  settings.gradle.kts
  gradlew / gradlew.bat          <- the wrapper (use this from now on)
  gradle/wrapper/                <- pins the Gradle version
  app/build.gradle.kts           <- the build file init wrote for you
  app/src/main/java/org/example/App.java
  app/src/test/java/org/example/AppTest.java
  \end{lstlisting}
  Two things to notice:
  \begin{itemize}
    \item the wrapper (\texttt{gradlew}) is part of the project now, so
      from here on every command starts with \texttt{./gradlew}
    \item tests already live where they belong:
      \texttt{src/test/java/...}, next to a sample JUnit 4 test
  \end{itemize}
\end{frame}

\begin{frame}[fragile]{Step 2b: reading the file \texttt{init} wrote}
  \texttt{app/build.gradle.kts}, trimmed to what matters:
  \begin{lstlisting}
plugins {
    application          // brings the run task (and java tasks)
}

repositories {
    mavenCentral()       // where libraries are fetched from
}

dependencies {
    testImplementation(libs.junit)  // JUnit 4, version lives
                                    // in gradle/libs.versions.toml
}

application {
    mainClass = "org.example.App"   // what ./gradlew run starts
}
  \end{lstlisting}
  Kotlin DSL notes: function-call style with double quotes
  (\texttt{mavenCentral()}), and versions kept in one shared catalog
  file instead of being scattered through the build.
\end{frame}

\begin{frame}[fragile]{Step 3: tasks come free with the plugins}
  You never define \texttt{build}, \texttt{test}, \texttt{clean} or
  \texttt{run}. The generated build file applies the right plugins, and
  the plugins add the tasks:
  \begin{itemize}
    \item the \texttt{java} plugin adds \texttt{build}, \texttt{test},
      \texttt{clean}, \texttt{check}, \texttt{jar}, \ldots
    \item the \texttt{application} plugin adds \texttt{run}
  \end{itemize}
  Prove it to yourself, from the project root:
  \begin{lstlisting}
./gradlew tasks --all
  \end{lstlisting}
\end{frame}

\begin{frame}[fragile]{Step 4: the terminal workflow (1 of 2)}
  From the project root:
  \begin{lstlisting}
./gradlew build   # compile + test + jar, the full check
./gradlew test    # only the tests
./gradlew clean   # delete the build/ folder
  \end{lstlisting}
  A healthy \texttt{build} ends with:
  \begin{lstlisting}
> Task :compileJava
> Task :test
> Task :build
BUILD SUCCESSFUL
  \end{lstlisting}
\end{frame}

\begin{frame}[fragile]{Step 4: the terminal workflow (2 of 2)}
  \begin{lstlisting}
./gradlew run               # run the app, prints Hello World!
./gradlew test --rerun-tasks  # run tests again even if
                              # nothing changed
./gradlew test --info         # extra detail when hunting
                              # a failure
  \end{lstlisting}
  Reports land under \texttt{build/} (git-ignored), for example
  \texttt{app/build/reports/tests/test/index.html}.
\end{frame}

\begin{frame}[fragile]{Adding production code: the rules}
  \texttt{init} gave you a working sample. When you add your own classes:
  \begin{enumerate}
    \item Put the file under \texttt{app/src/main/java/}, in the folder
      that matches its \texttt{package}.
    \item One public class per file; the file name matches the class name.
    \item Only code with a \texttt{main()} method can be \emph{run}.
      Helper classes are used through tests or through the app, never
      run directly.
  \end{enumerate}
  \begin{lstlisting}[language=Java]
package org.example;

public class Calculator {
    public int add(int a, int b) {
        return a + b;
    }

    public int divide(int a, int b) {
        if (b == 0) {
            throw new IllegalArgumentException("divisor must not be 0");
        }
        return a / b;
    }
}
  \end{lstlisting}
\end{frame}

\begin{frame}[fragile]{Adding tests: the rules (JUnit 4)}
  \begin{enumerate}
    \item Put the file under \texttt{app/src/test/java/}, in the same
      package folders as the code it tests.
    \item Name it after the class under test plus \texttt{Test}
      (\texttt{Calculator} becomes \texttt{CalculatorTest}).
    \item JUnit 4 shape: import \texttt{org.junit.Test}, annotate each
      test method with \texttt{@Test}, assert with
      \texttt{org.junit.Assert}. Test methods are \texttt{public void}.
  \end{enumerate}
  \begin{lstlisting}[language=Java,basicstyle=\ttfamily\tiny]
package org.example;

import static org.junit.Assert.assertEquals;
import org.junit.Test;

public class CalculatorTest {

    private final Calculator calc = new Calculator();

    @Test
    public void addsTwoNumbers() {
        assertEquals(5, calc.add(2, 3));
    }

    @Test(expected = IllegalArgumentException.class)
    public void divideByZeroThrows() {
        calc.divide(1, 0);
    }
}
  \end{lstlisting}
\end{frame}

\begin{frame}{References: the original Gradle docs}
  Everything in this deck comes from these pages. When in doubt, read
  the docs before asking:
  \begin{itemize}
    \item \href{https://docs.gradle.org/current/userguide/build_init_plugin.html}{Build Init plugin}: every \texttt{gradle init} prompt explained
    \item \href{https://docs.gradle.org/current/userguide/gradle_wrapper.html}{Gradle Wrapper}: how the pinned version works
    \item \href{https://docs.gradle.org/current/userguide/building_java_projects.html}{Building Java projects}: plugins, source layout, tasks
    \item \href{https://docs.gradle.org/current/userguide/java_testing.html}{Testing in Gradle}: how \texttt{test} finds and runs JUnit
    \item \href{https://docs.gradle.org/current/userguide/part1_gradle_init.html}{Beginner tutorial: Initializing the Project}: the full init-generated Java example
  \end{itemize}
\end{frame}

\begin{frame}[fragile]{Cheat sheet (take a photo)}
  \begin{lstlisting}
mkdir gradle-demo && cd gradle-demo
gradle init       # 1 Application, 1 Java, 21,
                  # 1 Single, 1 Kotlin, 1 JUnit 4

./gradlew build   # full check
./gradlew test    # tests only
./gradlew run     # run the app
./gradlew clean   # start fresh
  \end{lstlisting}
  Stuck? Ask in this order: read the terminal output, open the HTML report
  under \texttt{app/build/reports/}, then ask the TA. Thank you!
\end{frame}

\end{document}
